\documentclass[10pt,a4paper]{amsart}
\usepackage{amsmath,amssymb,amsthm,graphicx}
% Retained minimal preamble for the Stacks Project chapter more-morphisms.tex (More on Morphisms).
% Verbatim excerpt of the chapter preamble: Xy-pic setup (native preamble lines 14-18) and the native
% theorem declarations with the mathematical macros (native preamble lines 154-217). Copied mechanically;
% nothing added, removed or redefined. The wrapper class and the amsmath/amssymb/amsthm/graphicx base
% packages are supplied by the assembler, so the class/externaldocument part of the native preamble is omitted.
% For commutative diagrams we use Xy-pic
\usepackage[all]{xy}

% We use 2cell for 2-commutative diagrams.
\xyoption{2cell}
% Theorem environments.
%
\theoremstyle{plain}
\newtheorem{theorem}[subsection]{Theorem}
\newtheorem{proposition}[subsection]{Proposition}
\newtheorem{lemma}[subsection]{Lemma}

\theoremstyle{definition}
\newtheorem{definition}[subsection]{Definition}
\newtheorem{example}[subsection]{Example}
\newtheorem{exercise}[subsection]{Exercise}
\newtheorem{situation}[subsection]{Situation}

\theoremstyle{remark}
\newtheorem{remark}[subsection]{Remark}
\newtheorem{remarks}[subsection]{Remarks}

\numberwithin{equation}{subsection}

% Macros
%
\def\lim{\mathop{\mathrm{lim}}\nolimits}
\def\colim{\mathop{\mathrm{colim}}\nolimits}
\def\Spec{\mathop{\mathrm{Spec}}}
\def\Hom{\mathop{\mathrm{Hom}}\nolimits}
\def\Ext{\mathop{\mathrm{Ext}}\nolimits}
\def\SheafHom{\mathop{\mathcal{H}\!\mathit{om}}\nolimits}
\def\SheafExt{\mathop{\mathcal{E}\!\mathit{xt}}\nolimits}
\def\Sch{\mathit{Sch}}
\def\Mor{\mathop{\mathrm{Mor}}\nolimits}
\def\Ob{\mathop{\mathrm{Ob}}\nolimits}
\def\Sh{\mathop{\mathit{Sh}}\nolimits}
\def\NL{\mathop{N\!L}\nolimits}
\def\CH{\mathop{\mathrm{CH}}\nolimits}
\def\proetale{{pro\text{-}\acute{e}tale}}
\def\etale{{\acute{e}tale}}
\def\QCoh{\mathit{QCoh}}
\def\Ker{\mathop{\mathrm{Ker}}}
\def\Im{\mathop{\mathrm{Im}}}
\def\Coker{\mathop{\mathrm{Coker}}}
\def\Coim{\mathop{\mathrm{Coim}}}

% Boxtimes
%
\DeclareMathSymbol{\boxtimes}{\mathbin}{AMSa}{"02}

%
% Macros for moduli stacks/spaces
%
\def\QCohstack{\mathcal{QC}\!\mathit{oh}}
\def\Cohstack{\mathcal{C}\!\mathit{oh}}
\def\Spacesstack{\mathcal{S}\!\mathit{paces}}
\def\Quotfunctor{\mathrm{Quot}}
\def\Hilbfunctor{\mathrm{Hilb}}
\def\Curvesstack{\mathcal{C}\!\mathit{urves}}
\def\Polarizedstack{\mathcal{P}\!\mathit{olarized}}
\def\Complexesstack{\mathcal{C}\!\mathit{omplexes}}
% \Pic is the operator that assigns to X its picard group, usage \Pic(X)
% \Picardstack_{X/B} denotes the Picard stack of X over B
% \Picardfunctor_{X/B} denotes the Picard functor of X over B
\def\Pic{\mathop{\mathrm{Pic}}\nolimits}
\def\Picardstack{\mathcal{P}\!\mathit{ic}}
\def\Picardfunctor{\mathrm{Pic}}
\def\Deformationcategory{\mathcal{D}\!\mathit{ef}}
\usepackage{cleveref}
\crefname{originalreference}{original reference}{original references}
\Crefname{originalreference}{Original reference}{Original references}
\makeatletter
\newlabel{limits-lemma-ample-profinite-set-in-principal-affine}{{0E21}{source}{}{sourcepointer-1}{}}
\newlabel{limits-lemma-ample-profinite-set-in-principal-affine@cref}{{[originalreference][0][]0E21}{[1][0][]source}{}{}{}}
\makeatother
\crefname{originalreference}{original reference}{original references}
\Crefname{originalreference}{Original reference}{Original references}
\makeatletter
\newlabel{morphisms-lemma-integral-universally-closed}{{01WM}{source}{}{sourcepointer-2}{}}
\newlabel{morphisms-lemma-integral-universally-closed@cref}{{[originalreference][0][]01WM}{[1][0][]source}{}{}{}}
\makeatother
\crefname{originalreference}{original reference}{original references}
\Crefname{originalreference}{Original reference}{Original references}
\makeatletter
\newlabel{situation-pushout-along-closed-immersion-and-integral}{{0ECI}{source}{}{sourcepointer-3}{}}
\newlabel{situation-pushout-along-closed-immersion-and-integral@cref}{{[originalreference][0][]0ECI}{[1][0][]source}{}{}{}}
\makeatother
\title{Preparing a pushout square by matching affine opens along the closed piece}
\author{The Stacks Project Authors}
\date{stacks-project:more-morphisms.tex:local-snapshot-20261011; GNU FDL 1.2 or later}
\begin{document}
\maketitle
\noindent\textit{Editorial note.} Counted span is the author's own proof environment, source lines 19726-19755, immediately following the statement at 19719-19724.\par\medskip
\noindent\textit{Editorial note.} Two cross-chapter pointers (Limits, Morphisms) plus the chapter's own situation are recorded in assembly.external\_references.\par\medskip
\noindent\textit{Editorial note.} The author's phrase 'Now we (in some sense) repeat the argument' is retained verbatim.\par\medskip
\noindent\textit{Allowed standard prerequisites (editorial).}\begin{itemize}
\item The pushout situation of the section. chapter situation situation-pushout-along-closed-immersion-and-integral, cited in the statement. Fixes the closed immersion i : Z -> X and integral j : Z -> Y whose pushout is being prepared.
\item Integral morphisms are universally closed. Morphisms, lemma morphisms-lemma-integral-universally-closed, cited at proof line 19731. Makes j : Z -> Y closed, which is what lets an affine neighbourhood of y inside i\textasciicircum{}\{-1\}(U) be found.
\item Ample profinite sets in a principal affine open. Limits, lemma limits-lemma-ample-profinite-set-in-principal-affine, cited in the proof. Supplies the principal open D(g) inside the previous affine chart used in the second pass of the argument.
\end{itemize}\medskip
\ifdefined\hypertarget\hypertarget{sourcepointer-1}{}\fi
\noindent\textit{Source reference.} Original reference 0E21: focused-seeds/geometry/S9-morph-prepare-pushout-closed-integral/ref-tags.txt, source lines 4--4. Navigation pointer only: the counted proof cites this statement by \textbackslash{}ref as a standing input; the deduction the counted proof has to make is carried out inside the counted span, and no part of that deduction is delegated to this pointer.\par\medskip
\ifdefined\hypertarget\hypertarget{sourcepointer-2}{}\fi
\noindent\textit{Source reference.} Original reference 01WM: focused-seeds/geometry/S9-morph-prepare-pushout-closed-integral/ref-tags.txt, source lines 5--5. Navigation pointer only: the counted proof cites this statement by \textbackslash{}ref as a standing input; the deduction the counted proof has to make is carried out inside the counted span, and no part of that deduction is delegated to this pointer.\par\medskip
\ifdefined\hypertarget\hypertarget{sourcepointer-3}{}\fi
\noindent\textit{Source reference.} Original reference 0ECI: sources/stacks-more-morphisms/pinned-original-source.tex, source lines 19708--19708. Navigation pointer only: the counted proof cites this statement by \textbackslash{}ref as a standing input; the deduction the counted proof has to make is carried out inside the counted span, and no part of that deduction is delegated to this pointer.\par\medskip
\begin{lemma}
\label{lemma-prepare-pushout-along-closed-immersion-and-integral}
In Situation \ref{situation-pushout-along-closed-immersion-and-integral}
then for $y \in Y$ there exist affine opens $U \subset X$ and
$V \subset Y$ with $i^{-1}(U) = j^{-1}(V)$ and $y \in V$.
\end{lemma}

\begin{proof}
Let $y \in Y$. Choose an affine open $U \subset X$
such that $j^{-1}(\{y\}) \subset i^{-1}(U)$ (possible by assumption).
Choose an affine open $V \subset Y$ neighbourhood of $y$
such that $j^{-1}(V) \subset i^{-1}(U)$.
This is possible because $j : Z \to Y$ is a closed morphism
(Morphisms, Lemma \ref{morphisms-lemma-integral-universally-closed}) and
$i^{-1}(U)$ contains the fibre over $y$.
Since $j$ is integral, the scheme theoretic fibre $Z_y$
is the spectrum of an algebra integral over a field.
By Limits, Lemma \ref{limits-lemma-ample-profinite-set-in-principal-affine}
we can find an $\overline{f} \in \Gamma(i^{-1}(U), \mathcal{O}_{i^{-1}(U)})$
such that $Z_y \subset D(\overline{f}) \subset j^{-1}(V)$.
Since $i|_{i^{-1}(U)} : i^{-1}(U) \to U$ is a closed immersion
of affines, we can choose an $f \in \Gamma(U, \mathcal{O}_U)$
whose restriction to $i^{-1}(U)$ is $\overline{f}$.
After replacing $U$ by the principal open $D(f) \subset U$
we find affine opens $y \in V \subset Y$ and $U \subset X$ with
$$
j^{-1}(\{y\}) \subset i^{-1}(U) \subset j^{-1}(V)
$$
Now we (in some sense) repeat the argument. Namely, we choose
$g \in \Gamma(V, \mathcal{O}_V)$ such that $y \in D(g)$ and
$j^{-1}(D(g)) \subset i^{-1}(U)$ (possible by the same argument
as above). Then we can pick $f \in \Gamma(U, \mathcal{O}_U)$
whose restriction to $i^{-1}(U)$ is the pullback of $g$
by $i^{-1}(U) \to V$ (again possible by the same reason as above).
Then we finally have affine opens $y \in V' = D(g) \subset V \subset Y$
and $U' = D(f) \subset U \subset X$ with $j^{-1}(V') = i^{-1}(V')$.
\end{proof}

\end{document}
